% ---------------------------------------------------------------------------
% MSC PROPOSAL V1 — Cost-Aware Repository Change Localization with Sparse
% Impact Planning and Bounded Verification
%
% Supervisor-ready, conservative academic format. Content is transplantable
% into an official FCAI template later. No institutional form fields are
% fabricated. References verified against the living systematic review.
% ---------------------------------------------------------------------------
\documentclass[11pt,a4paper]{article}
\usepackage[T1]{fontenc}
\usepackage[utf8]{inputenc}
\usepackage{lmodern}
\usepackage{microtype}
\usepackage{amsmath,amssymb}
\usepackage{booktabs,tabularx,array}
\usepackage{enumitem}
\usepackage{graphicx}
\usepackage[margin=2.5cm]{geometry}
\usepackage{xcolor}
\usepackage{url}
\usepackage[colorlinks=true,linkcolor=blue,citecolor=blue,urlcolor=blue]{hyperref}
\emergencystretch=1.5em

\title{\textbf{Cost-Aware Repository Change Localization with Sparse Impact\\[2pt]
Planning and Bounded Verification}\\[6pt]
\large MSc Research Proposal (V1, supervisor-ready draft)}
\author{Ahmed Ehab}
\date{\today}

\begin{document}
\maketitle

\begin{abstract}
Large language models (LLMs) that edit a software repository must first decide
which files a requested change affects. This decision is costly: an explicit
file-by-file plan over a repository-scale candidate universe spends hundreds of
serialized records per call, and a second-stage reasoner that revisits the
scope can multiply inference cost by orders of magnitude. This thesis
investigates a representation of impact decisions that is cheap by construction
--- \emph{preserve-by-omission} sparse plans, which emit only non-preserve
decisions --- and asks when additional repository reasoning is worth its
inference cost. We frame the problem as a correctness--cost trade-off under
limited inference budgets and propose to combine (i) a sparse explicit
file-level impact policy, (ii) fair cheap/classical/LLM baselines under a
common file-selection protocol, and (iii) a \emph{bounded} second-stage
verification mechanism that revisits either the task or individual suspicious
candidates under a hard budget. Preliminary controlled and real-commit studies
show large serialization-cost reductions and meaningful zero-LLM lexical
signal, but also that task-level omission risk is not reliably predictable on
current development evidence. The thesis therefore evaluates the representation
and the correctness--cost Pareto frontier, and tests a pre-registered
contingency --- candidate-level bounded verification --- rather than requiring
a task-level risk scorer to succeed.
\end{abstract}

\section{Background and motivation}
\label{sec:background}

Repository-level coding agents and impact planners must select affected files
before editing. Two separable problems motivate this thesis. First,
\emph{impact inference}: did the model identify the right files? Second,
\emph{impact-policy representation}: can a complete file-action policy be
expressed without spending output budget on hundreds of repeated
\emph{preserve} decisions? Our controlled studies show that an explicit full
policy at a 16K token cap serializes on average 144 decision records per
candidate universe, while a sparse policy that emits only non-preserve
decisions serializes a fraction of that, with a large associated reduction in
completion tokens and cost.

The second motivating fact is that second-stage reasoning is expensive.
A graph-guided baseline under a common protocol consumed on the order of a
million tokens per task on the same held-out tasks where our sparse first-pass
stays below a few thousand tokens. If a cheap first pass already localizes the
bulk of affected files, then the remaining question is whether a \emph{bounded}
second look can recover missed files without paying the full always-on cost.

\section{Problem statement}
\label{sec:problem}

Given a repository revision $P$ (parent commit), a natural-language change
request $q$, and a candidate universe $U(P)$ of production files visible at
$P$, a localization system produces an affected-file set $\hat{S}(P,q)$.
The hidden observed change set $S^*(P)$ is used for evaluation only. The thesis
studies correctness---cost trade-offs under explicit budgets:

\begin{itemize}
  \item \textbf{Representation cost.} How do sparse (\emph{preserve-by-
  omission}) and explicit impact policies compare in output tokens, serialized
  records, cost, latency, and selection fidelity?
  \item \textbf{Bounded second-stage verification.} Can a second mechanism that
  revisits either the task (\emph{Route A}) or individual suspicious omitted
  candidates (\emph{Route B}) recover false negatives more efficiently than
  always-on or random matched-budget verification?
\end{itemize}

The observed historical diff is treated as an \emph{observed change-set proxy},
never as semantic ground truth.

\section{Gap from existing work}
\label{sec:gap}

Existing lines --- agentic graph-guided localization (LocAgent,
GraphLocator), repository memory / retrieval-conditioned generation (RepoCoder,
RepoAgent), explicit edit planning (CodePlan), agentless localize-then-repair
(Agentless), and classical change-impact/propagation analysis --- either
target semantic fidelity without a cost discipline, or address cheap lexical
localization without a second-stage omission-recovery mechanism. None combines
an explicit sparse impact-policy representation, a fair cheap/classical/LLM
comparison under a common file-selection protocol, and a budgeted
omission-aware verifier. Our contribution is not any single component but the
combination and its honest, cost-aware evaluation on real commits.

\section{Research objectives}
\label{sec:objectives}

\begin{enumerate}
  \item Build reproducible real-commit localization datasets and protocols
        (djangoCMS V2-LARGE with an untouched internal test; Saleor as a second
        real system; NestJS as a cross-ecosystem check).
  \item Quantify the representation cost of sparse explicit impact policies and
        their file-selection fidelity under shared protocols.
  \item Evaluate cheap, classical, and LLM localization under common
        file-level metrics and explicit budgets.
  \item Design and test a bounded second-stage verification mechanism
        (Route A task-level routing, or Route B candidate-level verification),
        pre-registered and evaluated under matched budgets.
  \item Produce correctness--cost Pareto evidence per repository, with
        macro-averaged cross-repository summaries and a repository-identity
        confound check.
\end{enumerate}

\section{Research questions}
\label{sec:rq}

\begin{description}
  \item[RQ1.] How does sparse explicit impact-policy representation affect
  inference/output cost and file-selection fidelity?
  \item[RQ2.] How accurately can cheap lexical/static methods localize impacted
  files relative to LLM planners under shared protocols?
  \item[RQ3.] Can bounded second-stage verification recover missed impacted
  files more efficiently than always-on or random matched-budget verification?
  \item[RQ4.] How do these correctness--cost trade-offs transfer across
  repositories and ecosystems?
\end{description}

If task-level risk modeling remains unsupported by evidence, RQ3 is answered
by the candidate-level (Route B) mechanism; \emph{RiskScorer success is not a
thesis requirement}.

\section{Current preliminary evidence}
\label{sec:evidence}

All numbers are from our audited studies (this repository); they are
development or controlled evidence and are reported with their bounds.

\begin{itemize}
  \item \textbf{Preserve-by-omission representation effect (controlled).}
  In a 16K-cap controlled study (M1B, 60 cells, 6 curated tasks), Sparse-v2
  reduced mean completion tokens by $\sim$90\%, mean serialized decision
  records by $\sim$95\% (corrected metric), recorded API cost by $\sim$82\%,
  and latency by $\sim$63\% relative to an explicit Full-v2 policy, with both
  arms fully valid. No universal semantic-superiority claim is made.
  \item \textbf{Real-commit P1} (60 cells; 10 held-out real djangoCMS commits;
  2 arms, 3 nested repetitions): semantic selection is difficult --- FNR is
  high; there is \emph{no} Full-vs-Sparse superiority claim (paired deltas'
  confidence intervals cross zero). Sparse execution is substantially cheaper.
  \item \textbf{LocAgent (P5)}: system-level shared-protocol context only ---
  402 calls / 32.8M tokens / $\sim$9.93 USD estimated on 10 tasks with a 50\%
  empty localization rate; not a reproduction of the published fine-tuned
  result.
  \item \textbf{Cheap baselines (Protocol A)}: BM25 provides a meaningful
  zero-LLM localization signal on development data; graph\@K $\approx$
  path\_token\@K with lexical seeds; K is an operating-point curve, not a
  fixed configuration.
  \item \textbf{Sparse-v2 development inference (90 cells, TRAIN/VALIDATION)}:
  live API, 90/90 valid, 490,747 tokens / \$\ 0.184 recorded cost; Sparse
  omission prevalence 26/30. A task-level RiskScorer is \emph{not} justified on
  this evidence (4 negatives $<$ 10; no reliable feature above the random
  band).
  \item \textbf{V2 development inference (2026-09-16)}: 431/450 cells within
  the authorized ceilings (fail-closed token-ceiling stop at 2.5M tokens,
  \$\ 0.873); 144 new V2 development tasks; the v1 candidate omission signal
  does not replicate within V2 (apparent pooled signal is a candidate-universe
  size artifact). No multivariable RiskScorer is fitted.
  \item \textbf{Classical/static CIA baseline}: implemented on 150 V2
  development cases (zero LLM); BM25 and graph-ranked seeds are the strongest
  arms; reverse-dependency closure arms are weaker.
\end{itemize}

These results \emph{motivate} a larger development dataset (V2-LARGE),
protection of an untouched internal test, and the bounded-verification design;
they are not rewritten as a success.

\section{Proposed architecture}
\label{sec:architecture}

\begin{enumerate}
  \item \textbf{Sparse first-pass impact plan}: explicit file-level policy with
  preserve-by-omission serialization; cheap deterministic and LLM planner
  variants under the same contract.
  \item \textbf{Omission-risk estimation (contingency)}: only if a larger
  development set establishes a reliable, stable signal; a simple rule /
  logistic model, no neural model, no broad feature search.
  \item \textbf{Second-stage verification}:
    \begin{itemize}
      \item \emph{Route A}: task-level selective routing (escalate the whole
      task when risk is high).
      \item \emph{Route B}: candidate-level bounded verification --- identify
      suspicious omitted candidates using independent structural/history/
      retrieval evidence, verify only those under a hard budget, compare with
      always-verify and random matched-budget verification.
    \end{itemize}
\end{enumerate}

\section{Dataset methodology}
\label{sec:dataset}

\begin{itemize}
  \item \textbf{Real history, parent-only inputs.} Inference uses the parent
  commit $P$; the target $T$ is hidden. Live repository HEAD is never used
  (it moves and can leak future state).
  \item \textbf{Observed change-set proxy.} $P\!\to\!T$ changed production
  files are an observed proxy, never semantic gold.
  \item \textbf{V2-LARGE (Stage 1).} djangoCMS eligible pool (329 independent
  candidates after frozen filters/deduplication; 289 untouched beyond the
  exposed v1 40) split into DEV, an untouched internal test, and a reserve.
  The split is metadata-only and frozen before any model result.
  \item \textbf{Saleor (Stage 2)} and \textbf{NestJS (Stage 3)}: suitability
  audited (SUITABLE-WITH-DEVIATIONS); protocols frozen; inference is deferred
  until their sampling frames are reconstructed and splits frozen.
\end{itemize}

\section{Baselines}
\label{sec:baselines}

Random/cheap control, BM25, adaptive cheap retrieval (if justified),
classical/static CIA (dependency-propagation closure), Sparse planner, Full
planner, faithful LocAgent/reference agent (when reproducible), Sparse +
Always Verify, Sparse + Random Verify (matched budget), Sparse + Selective/
Bounded Verify, and an Oracle-routing upper bound. If task-level routing is
unsupported, the selective arm becomes candidate-level bounded verification.

\section{Evaluation dimensions}
\label{sec:eval}

Precision, recall, F1, FNR, FN recovery; selected-candidate count;
escalation/verification rate; calls; tokens; cost; latency; failures;
graph/index build cost separated from query cost; and the correctness--cost
Pareto frontier. Per-repository metrics are primary; macro-averages across
repositories are reported; pooled metrics are secondary; a repository-identity
confound check is included.

\section{Statistical plan}
\label{sec:stats}

\begin{itemize}
  \item Task is the independent unit; nested repetitions are never treated as
        independent observations.
  \item Pre-registered class-balance gate: risk modeling requires at least 10
        independent positives and 10 negatives, a predeclared feature above the
        random band, and direction stability across development/validation.
  \item Bootstrap confidence intervals over tasks; random-risk bands for
        AUROC at the observed $n$; multiple-comparison awareness (features
        above the band are compared with the chance expectation).
  \item Sample-size planning uses development evidence only (never test set);
        the number of tasks is chosen from the frozen sampling frame, not a
        round number.
\end{itemize}

\section{Cross-repository generalization plan}
\label{sec:gen}

Stage 1 djangoCMS V2-LARGE (within-repo large evidence + untouched internal
test) $\to$ Stage 2 Saleor (different real system, same Python/Django stack)
$\to$ Stage 3 NestJS (cross-language/cross-framework evidence). Each stage
states its claim boundary: within-repo development + internal confirmation;
cross-repository confirmation; cross-ecosystem evidence. NestJS is preferred
but not forced; a predeclared eligible-pool yield criterion and a JabRef
backup are frozen.

\section{Threats to validity}
\label{sec:threats}

\begin{itemize}
  \item \textbf{Internal:} the observed proxy is not semantic gold; leakage
  control (allow-intent-path-leakage off, parent-only inputs) is mandatory;
  no hidden-test peeking.
  \item \textbf{External:} $n$ and single-repository generalization; V2 and
  Saleor/NestJS are the mitigation, with per-repository primary metrics.
  \item \textbf{Statistical:} low negatives in the development label;
  multiple comparisons; wide random bands --- addressed by pre-registration,
  class-balance gates, and bootstrap uncertainty.
  \item \textbf{Representation:} LLM serialization cost spans providers; we
  report both recorded and frozen-pricing cost and keep provider pinned as far
  as allowed.
\end{itemize}

\section{Expected contributions}
\label{sec:contrib}

\begin{enumerate}
  \item Sparse explicit impact-plan representation with controlled evidence.
  \item A reproducible real-commit impact-localization dataset and protocol
        (repository, eligibility, split, proxy, hashes).
  \item Fair cheap/classical/LLM/agent comparison under common file-level
        metrics and explicit budgets.
  \item A bounded/selective-compute verification protocol (Route A/B) if
        supported by evidence.
  \item Cross-repository correctness--cost evidence.
  \item Reproducibility artifacts and leakage-control methodology.
\end{enumerate}

No positive algorithmic result is promised.

\section{Work plan / timeline}
\label{sec:timetable}

\begin{tabular}{ll}
\toprule
\textbf{Window} & \textbf{Activity} \\
\midrule
2026-09/10 & Proposal V1; dataset V2-LARGE freeze + build \\
2026-10 & V2 development inference; classical + cheap baselines; Route B design \\
2026-11 & Bounded verification experiments; Saleor protocol + inference \\
2026-12 & NestJS cross-ecosystem check; Pareto analysis \\
2027-01/02 & Analysis, replication, threat revisions \\
2027-03/04 & Write-up; internal test confirmation \\
2027-05 & Submission-ready thesis \\
\bottomrule
\end{tabular}

\section{References}
\label{sec:refs}

The authoritative bibliography is \texttt{references.bib}. Verified primary
sources include (venue where confirmed): Agentless (arXiv, 2024); CodePlan
(ICSE 2024); RepoCoder (EMNLP 2023); AutoCodeRover (ISSTA 2024); RepoGraph
(ICLR 2025); LocAgent (arXiv 2025, upstream pin evaluated under a shared
protocol); the Shichao Zhang adaptive-kNN/cost-sensitive line (TKDE 2021/2022;
Neurocomputing 2020); GraphLocator (FSE 2026, LLM causal-issue-graph
localization); and the classical change-impact/ripple-effect line. Preprints
are marked as such in the bibliography.

\end{document}